\chapter{4x10\textsuperscript{9}D.C.}\label{4x10^9-8}

\hfill\break

Fuori da Padang passammo dall'ambulanza di Nijon a una macchina privata
con un autista Minangkabau, che fece scendere me, Ibu Ina ed En in un
complesso per trasporto merci sull'autostrada costiera. Cinque magazzini
enormi col tetto di lamiera si estendevano su una piana di ghiaia scura
tra ammassi conici di cemento sfuso coperti da teloni e una cisterna
ferroviaria corrosa e inattiva su un binario morto. L'ufficio principale
era un edificio basso in legno con un'insegna che recitava B{AYUR}
F{ORWARDING}, ovvero Spedizioni Bayur.

Bayur Forwarding, disse Ina, era una delle attività del suo ex marito
Jala e fu proprio Jala che ci venne incontro nella stanza della
reception. Era un uomo robusto dalle guance rosse come mele con un
completo giallo canarino - sembrava una caraffa Toby vestita per i
tropici. Lui e Ina si abbracciarono alla maniera dei divorziati in buoni
rapporti, poi Jala mi strinse la mano e si chinò a stringere quella di
En. Mi presentò alla sua receptionist come un importatore di olio di
palma di Suffolk, forse per evitare rischi nel caso in cui venisse
interrogata dalla nuova Reformasi. Poi ci accompagnò alla sua BMW a
celle a combustibile di sette anni e ci dirigemmo a sud verso Teluk
Bayur, Jala e Ina seduti davanti, io ed En dietro.

Teluk Bayur, il grande porto in acque profonde a sud della città di
Padang, era la fonte di tutti i guadagni di Jala. Trent'anni prima,
disse, Teluk Bayur era stato un bacino sonnolento di fango sabbioso con
servizi portuali modesti e un commercio prevedibile di carbone, olio di
palma non raffinato e fertilizzanti. In quel periodo, invece, grazie al
boom economico della Restaurazione Nagari e all'esplosione demografica
dalla comparsa dell'Arco, Teluk Bayur era diventato un bacino portuale
all'avanguardia con banchine e ormeggi di prim'ordine, un enorme
complesso di stoccaggio e così tanti comfort moderni che persino Jala
alla fine perse l'interesse nel tenere il conto di tutti i
rimorchiatori, i capannoni, le gru e i carrelli elevatori da container.
«Jala è fiero di Teluk Bayur» disse Ina. «Non c'è alto funzionario o
quasi che non abbia corrotto.»

«Nessuno più in alto degli Affari generali» la corresse Jala.

«Sei troppo modesto.»

«C'è qualcosa che non va nel fare soldi? Ho avuto troppo successo? È un
crimine diventare qualcuno?»

Ina inclinò la testa e disse: «Sono ovviamente domande retoriche.»

Chiesi se ci saremmo imbarcati direttamente a Teluk Bayur.

«No, non direttamente» disse Jala. «Vi porto in un posto sicuro sul
molo. Non è semplice come salire su un'imbarcazione e mettersi comodi.»

«Non ci sono navi?»

«Certo, ce n'è una. La \emph{Capetown Maru}, un bel mercantile non
troppo grande. Sta caricando caffè e spezie proprio ora. Quando la stiva
sarà piena, i debiti pagati e i permessi firmati, allora il carico umano
salirà a bordo. In modo discreto, mi auguro.»

«E Diane? Diane è a Teluk Bayur?»

«Presto» disse Ina, indirizzando a Jala un'occhiata significativa.

«Sì, presto» disse lui.

\separatorefiori

Un tempo Teluk Bayur poteva anche essere stato un porto commerciale
sonnolento, ma come qualsiasi porto moderno era diventato una città a sé
stante, una città fatta non per le persone ma per i carichi di merce. Il
porto in sé era delimitato e recintato, ma le attività secondarie vi
erano cresciute attorno come bordelli fuori da una base militare:
noleggiatori e trasportatori secondari, cooperative di camionisti
indipendenti che viaggiavano su TIR ricondizionati e gestivano depositi
colabrodo di carburante. Vi passammo attraverso velocemente, senza
fermarci. Jala voleva trovarci una sistemazione prima che calasse il
Sole.

La stessa baia di Bayur era un ferro di cavallo di acqua salata oleosa.
Pontili e moli vi lappavano dentro come lingue di cemento. A ridosso
della costa c'era il caos ordinato del commercio su larga scala,
magazzini di prima e seconda linea e stazioni di accatastamento delle
merci, gru come gigantesche mantidi che banchettavano con le stive di
navi portacontainer ancorate. Ci fermammo a un posto di guardia
presidiato lungo il confine di una recinzione in acciaio e dal
finestrino dell'auto Jala allungò qualcosa all'addetto alla sicurezza:
un permesso, una mazzetta o entrambe le cose. La guardia gli fece cenno
di passare, Jala salutò amichevolmente e proseguì all'interno, seguendo
una fila di cisterne di olio di palma non raffinato e benzina avio a una
velocità che sembrava sconsiderata. Disse: «Ho sistemato le cose in modo
che possiate passare qui la notte. Ho un ufficio in uno dei magazzini
del molo E. Non c'è niente se non sacchi di cemento, nessuno che possa
disturbarvi. Domattina porterò Diane Lawton.»

«E poi partiamo?»

«Sii paziente. Non siete gli unici a fare \emph{rantau}\ldots{} ma
quelli che danno più nell'occhio. Ci potrebbero essere complicazioni.»

«Del tipo?»

«Ovviamente la nuova Reformasi. Di tanto in tanto la polizia perlustra
le zone portuali, in cerca di clandestini o gente che vuole attraversare
l'Arco. Di solito ne trovano qualcuno. O più di qualcuno, a seconda di
chi è stato pagato. Al momento ci sono forti pressioni da Jakarta,
perciò chi lo sa? Si parla anche di un'azione sindacale. Il sindacato
dei portuali è davvero attivo. Salperemo prima che inizi qualsiasi
conflitto, se siamo fortunati. Perciò temo che questa notte dovrai
dormire sul pavimento e al buio, per ora porto Ina ed En dagli altri
abitanti del villaggio.»

«No» reagì Ina con fermezza. «Rimarrò con Tyler.»

Jala fece una pausa. Poi la guardò e disse qualcosa in minang.

«Non è divertente» replicò lei. «E non è vero.»

«Allora, cosa? Non ti fidi che riesca a tenerlo al sicuro?»

«Cos'ho mai ottenuto fidandomi di te?»

Jala fece un gran sorriso. Aveva i denti color tabacco. «Avventura.»

«Sì, proprio» disse Ina.

\separatorefiori

E così io e Ibu Ina finimmo all'estremità nord di un complesso di
magazzini lontano dai moli, in un'orrenda stanza rettangolare che era
stata l'ufficio dell'ispettore, disse Ina, finché l'edificio non venne
temporaneamente chiuso in attesa di riparare il tetto poroso.

Una parete della stanza era una finestra di vetro retinato. Guardai
dentro un deposito cavernoso imbiancato dalla polvere di cemento. Le
travi di sostegno in acciaio nascevano come costole arrugginite dal
pavimento coperto di fango e acqua stagnante.

L'unica luce proveniva dalle lampade di sicurezza disposte ad ampi
intervalli lungo le pareti. Degli insetti si erano intrufolati dai buchi
dell'edificio e volavano in nugoli attorno alle lampadine ingabbiate o
vi morivano ammucchiandosi sotto. Ina riuscì a far funzionare una
lampada da tavolo. In un angolo erano stati accatastati dei cartoni
vuoti. Aprii quelli più asciutti e li impilai per farne un paio di
materassi alla buona. Non c'erano coperte, ma era una notte calda.
Eravamo vicini alla stagione dei monsoni.

«Pensi di riuscire a dormire?» chiese Ina.

«Non è l'Hilton, ma non c'è di meglio.»

«No, non mi riferivo a quello. Volevo dire il rumore. Riesci a dormire
col rumore?»

Teluk Bayur di notte non chiudeva. Il carico e lo scarico andavano
avanti ventiquattro ore al giorno. Non potevamo vederlo ma riuscivamo a
sentirlo, il rumore di motori pesanti e metallo sotto pressione e il
fragore intermittente dei container in transito. «Ho dormito in
situazioni peggiori» dissi.

«Ne dubito,» replicò Ina, «ma è gentile da parte tua.»

Nessuno dei due riuscì a dormire per ore. E quindi ci sedemmo accanto al
bagliore della lampada da tavolo e ogni tanto parlavamo. Ina mi chiedeva
di Jason.

Le avevo permesso di leggere qualche lungo passaggio che avevo scritto
durante la mia malattia. La transizione di Jason a Quarto, disse,
sembrava essere stata meno complicata della mia. No, ribattei. Avevo
semplicemente omesso di aggiungere i particolari della \emph{padella}.

«E la sua memoria? Non c'è stata nessuna perdita? Era sereno?»

«Non ne ha parlato molto, ma sono certo che fosse preoccupato.»

Infatti, riemergendo da una delle febbri ricorrenti mi aveva chiesto di
documentare per lui la sua vita: ``Scrivilo per me, Ty'' aveva detto.
``Scrivilo nel caso io dimentichi.''

«Quindi per lui nessuna grafomania.»

«No. La grafomania si verifica quando il cervello comincia a ricreare le
sue facoltà verbali. È solo uno dei sintomi possibili. I suoni che
emetteva forse ne erano la sua manifestazione.»

«Tu l'hai imparato da Wun Ngo Wen.»

Sì, o dai suoi archivi medici che avevo studiato in seguito.

Ina era ancora affascinata da Wun Ngo Wen. «Di quell'avvertimento alle
Nazioni Unite, sulla sovrappopolazione e l'esaurimento delle risorse,
Wun ne ha mai discusso con te? Intendo nel periodo precedente a\ldots»

«Ho capito. Sì, un po'.»

«Cosa ti ha detto?»

Accadde in una delle nostre conversazioni sullo scopo ultimo degli
Ipotetici. Wun mi aveva disegnato un diagramma, che riprodussi per Ina
sul pavimento di parquet polveroso: una linea orizzontale e una
verticale che componevano un grafico. La linea verticale era la
popolazione, quella orizzontale il tempo. Una linea di tendenza
irregolare attraversava il grafico più o meno orizzontalmente.

«La popolazione nel tempo» disse Ina. «Lo capisco, ma cosa vogliamo
misurare, in effetti?»

«Una qualsiasi popolazione animale in un ecosistema relativamente
stabile. Potrebbero essere le volpi dell'Alaska o le scimmie urlanti del
Belize. La popolazione oscilla secondo fattori esterni, come un inverno
freddo o un aumento dei predatori, ma è stabile, quantomeno nel breve
periodo.»

E invece, aveva detto Wun, che cosa succede se consideriamo nel lungo
periodo una specie intelligente in grado di utilizzare strumenti?
Disegnai a Ina lo stesso grafico precedente, solo che stavolta la linea
di tendenza curvava costantemente verso la direzione verticale.

«Qui succede che la popolazione - possiamo dire semplicemente ``le
persone'' - le persone imparano a mettere insieme le loro abilità. Non
solo come si spacca una selce ma come insegnare ad altre persone a
spaccare le selci e come dividere il lavoro in termini economici. La
collaborazione crea più cibo. La popolazione cresce. Altre persone
collaborano in maniera più efficace e generano nuove competenze.
Agricoltura. Zootecnia. Lettura e scrittura, il che significa che le
abilità possono essere condivise in modo più efficace tra i vivi ma
anche ereditate da generazioni morte da tempo.»

«Quindi la curva sale in modo sempre più ripido,» disse Ina, «finché non
ci troveremo ad affogare in noi stessi.»

«Ah, ma questo non accade. Ci sono altre forze che lavorano per spostare
la curva verso destra. La prosperità crescente e l'acume tecnologico in
realtà giocano a nostro favore. Le persone ben nutrite e sicure tendono
a voler limitare la loro riproduzione. La tecnologia e una cultura
flessibile danno loro i mezzi per farlo. In definitiva, o almeno così ha
detto Wun, la curva ritenderà verso la posizione orizzontale.»

Ibu Ina sembrava confusa. «Quindi \emph{non c'è} nessun problema? Né
inedia, né sovrappopolazione?»

«Purtroppo, la linea per la popolazione della Terra è ancora ben lontana
dall'orizzontale. E andremo incontro a condizioni limitanti.»

«Condizioni limitanti?»

Un altro diagramma. Stavolta mostrava una linea di tendenza a forma di S
corsiva, piatta in cima. Su questa segnai due linee orizzontali
parallele: una ben al di sopra della linea di tendenza, contrassegnata
da ``A'', e una che l'attraversava nella curva superiore, contrassegnata
da ``B''.

«Cosa rappresentano queste linee?» chiese Ina.

«Rappresentano entrambe la sostenibilità del pianeta. La quantità di
terra coltivabile disponibile per l'agricoltura, il combustibile e le
materie prime per sostenere tecnologia, aria pulita e acqua. Il
diagramma mostra la differenza tra una specie intelligente vincente e
una non vincente. Una specie che raggiunge il suo apice prima del limite
ha un potenziale di sopravvivenza a lungo termine. Una specie vincente
può continuare a fare tutte quelle cose che i futuristi sognavano:
espandersi nel sistema solare o persino nella galassia, manipolare il
tempo e lo spazio.»

«Che esagerazione» disse Ina.

«Non criticare. L'alternativa è peggiore. Una specie che arriva ai
limiti della sostenibilità prima di stabilizzare la sua popolazione
probabilmente è condannata. Inedia di massa, fallimento tecnologico e un
pianeta così svuotato dal primo fiore della civiltà che gli mancano i
mezzi per ricostruirsi.»

«Capisco.» Rabbrividì. «Quindi quale dei due siamo noi? Caso A o Caso B?
Wun te l'ha detto?»

«Poteva dire con certezza solo che entrambi i pianeti, Terra e Marte,
stavano per arrivare al limite. E che gli Ipotetici sono intervenuti
prima che potesse accadere.»

«Ma \emph{perché} sono intervenuti? Cosa si aspettano da noi?»

Era una domanda alla quale il popolo di Wun non aveva trovato risposta.

E nemmeno noi.

No, non era del tutto vero. Jason Lawton aveva trovato una specie di
risposta.

Ma non ero ancora pronto a parlarne.

\separatorefiori

Ina sbadigliò e così rimossi i segni sul pavimento polveroso. Spense la
luce della lampada da tavolo. Le lampade da cantiere sparpagliate in
giro emanavano un bagliore esausto. Fuori dal magazzino, più o meno ogni
cinque secondi si udiva un suono simile al rintocco smorzato di una
campana enorme.

«Tic toc» sbuffò Ina, sistemandosi sul materasso di cartoni ammuffiti.
«Mi ricordo quando gli orologi ticchettavano, Tyler. E tu? Quegli
orologi all'antica?»

«Ce n'era uno nella cucina di mia madre.»

«Ci sono tantissimi generi di tempo. Il tempo con cui misuriamo la
nostra vita. Mesi e anni. O il \emph{grande} tempo, il tempo che innalza
montagne e crea stelle. O tutte le cose che accadono tra un battito di
cuore e l'altro. È difficile vivere in tutti quei generi di tempo. È
facile \emph{dimenticare} che si vive in ognuno di loro.»

Il clangore metronomico continuava.

«Sembri un Quarto» dissi.

Nella penombra riuscii a scorgere solo il suo sorriso stanco.

«Penso che una vita mi basti» disse lei.

\separatorefiori

Al mattino ci svegliammo al rumore di una porta a soffietto ripiegata
fino ai fermi, un'esplosione di luce, Jala che ci chiamava.

Mi precipitai giù per le scale. Jala aveva già percorso metà magazzino e
Diane lo seguiva, camminando lentamente.

Mi avvicinai e la chiamai per nome.

Cercò di sorridere, ma aveva i denti serrati e il viso stranamente
pallido. Avevo già notato che si premeva un panno arrotolato sopra
l'anca e che sia il panno che la camicetta di cotone erano di un colore
rosso vivo per il sangue fuoriuscito.